\documentclass[10pt,a4paper]{amsart}
\usepackage{amsmath,amssymb,amsthm,graphicx}
\usepackage{verbatim}
\newenvironment{reference}{\comment}{\endcomment}
%\newenvironment{reference}{}{}
\newenvironment{slogan}{\comment}{\endcomment}
\newenvironment{history}{\comment}{\endcomment}

% For commutative diagrams we use Xy-pic
\usepackage[all]{xy}

% We use 2cell for 2-commutative diagrams.
\xyoption{2cell}
% Theorem environments.
%
\theoremstyle{plain}
\newtheorem{theorem}[subsection]{Theorem}
\newtheorem{proposition}[subsection]{Proposition}
\newtheorem{lemma}[subsection]{Lemma}

\theoremstyle{definition}
\newtheorem{definition}[subsection]{Definition}
\newtheorem{example}[subsection]{Example}
\newtheorem{exercise}[subsection]{Exercise}
\newtheorem{situation}[subsection]{Situation}

\theoremstyle{remark}
\newtheorem{remark}[subsection]{Remark}
\newtheorem{remarks}[subsection]{Remarks}

\numberwithin{equation}{subsection}

% Macros
%
\def\lim{\mathop{\mathrm{lim}}\nolimits}
\def\colim{\mathop{\mathrm{colim}}\nolimits}
\def\Spec{\mathop{\mathrm{Spec}}}
\def\Hom{\mathop{\mathrm{Hom}}\nolimits}
\def\Ext{\mathop{\mathrm{Ext}}\nolimits}
\def\SheafHom{\mathop{\mathcal{H}\!\mathit{om}}\nolimits}
\def\SheafExt{\mathop{\mathcal{E}\!\mathit{xt}}\nolimits}
\def\Sch{\mathit{Sch}}
\def\Mor{\mathop{\mathrm{Mor}}\nolimits}
\def\Ob{\mathop{\mathrm{Ob}}\nolimits}
\def\Sh{\mathop{\mathit{Sh}}\nolimits}
\def\NL{\mathop{N\!L}\nolimits}
\def\CH{\mathop{\mathrm{CH}}\nolimits}
\def\proetale{{pro\text{-}\acute{e}tale}}
\def\etale{{\acute{e}tale}}
\def\QCoh{\mathit{QCoh}}
\def\Ker{\mathop{\mathrm{Ker}}}
\def\Im{\mathop{\mathrm{Im}}}
\def\Coker{\mathop{\mathrm{Coker}}}
\def\Coim{\mathop{\mathrm{Coim}}}

% Boxtimes
%
\DeclareMathSymbol{\boxtimes}{\mathbin}{AMSa}{"02}

%
% Macros for moduli stacks/spaces
%
\def\QCohstack{\mathcal{QC}\!\mathit{oh}}
\def\Cohstack{\mathcal{C}\!\mathit{oh}}
\def\Spacesstack{\mathcal{S}\!\mathit{paces}}
\def\Quotfunctor{\mathrm{Quot}}
\def\Hilbfunctor{\mathrm{Hilb}}
\def\Curvesstack{\mathcal{C}\!\mathit{urves}}
\def\Polarizedstack{\mathcal{P}\!\mathit{olarized}}
\def\Complexesstack{\mathcal{C}\!\mathit{omplexes}}
% \Pic is the operator that assigns to X its picard group, usage \Pic(X)
% \Picardstack_{X/B} denotes the Picard stack of X over B
% \Picardfunctor_{X/B} denotes the Picard functor of X over B
\def\Pic{\mathop{\mathrm{Pic}}\nolimits}
\def\Picardstack{\mathcal{P}\!\mathit{ic}}
\def\Picardfunctor{\mathrm{Pic}}
\def\Deformationcategory{\mathcal{D}\!\mathit{ef}}
\usepackage{cleveref}
\crefname{originallemma}{original lemma}{original lemmas}
\Crefname{originallemma}{Original lemma}{Original lemmas}
\makeatletter
\newlabel{lemma-colimit-sheaves}{{00WI}{source}{}{sourcepointer-1}{}}
\newlabel{lemma-colimit-sheaves@cref}{{[originallemma][0][]00WI}{[1][0][]source}{}{}{}}
\makeatother
\crefname{originallemma}{original lemma}{original lemmas}
\Crefname{originallemma}{Original lemma}{Original lemmas}
\makeatletter
\newlabel{lemma-covering-surjective-after-sheafification}{{00WT}{source}{}{sourcepointer-2}{}}
\newlabel{lemma-covering-surjective-after-sheafification@cref}{{[originallemma][0][]00WT}{[1][0][]source}{}{}{}}
\makeatother
\title{Quasi-compact objects and finite subcoverings of a surjection of sheaves}
\author{The Stacks Project Authors}
\date{stacks-project:sites.tex:local-snapshot-20261012; GNU FDL 1.2 or later}
\begin{document}
\maketitle
\noindent\textit{Editorial note.} All delivered mathematical text is the author's own, copied from the pinned Stacks Project chapter file sites-cohomology.tex; no statement or proof step was rewritten, reordered or removed.\par\medskip
\noindent\textit{Editorial note.} The proof is the author's complete argument; every step the author writes out is inside the counted span and no closing sentence is a placeholder.\par\medskip
\noindent\textit{Editorial note.} Every labelled reference occurring inside the counted spans is carried in assembly.external\_references as a non-load-bearing navigation pointer (same-chapter ones at their definition site, cross-chapter ones at the invocation inside the counted span), so no load-bearing mathematics is replaced by a stub. No bibliography citation occurs inside the counted spans.\par\medskip
\noindent\textit{Editorial note.} The checked chapter list for same-name collisions is the set of stacks chapters staged in this repository: chow, cohomology, divisors, etale-cohomology, flat, more-algebra, more-morphisms, morphisms, spaces-more-morphisms, varieties together with this card's chapter sites-cohomology.\par\medskip
\noindent\textit{Editorial note.} The preamble is the chapter's own retained minimal preamble (sources/stacks-sites-cohomology/retained-preamble.tex), a verbatim excerpt of the native preamble: the Xy-pic setup and the native theorem declarations with the mathematical macros.\par\medskip
\noindent\textit{Allowed standard prerequisites (editorial).}\begin{itemize}
\item Coproducts of sheaves as sheafified coproducts. Same chapter, Lemma lemma-colimit-sheaves. Describes the sections of a coproduct of sheaves through a covering, which is used to replace a section by one coming from a single summand.
\item Quasi-compact objects of a topos. Same chapter, the definition of a quasi-compact object and its characterization by coverings. Supplies the property of a sieve or covering used in both directions.
\end{itemize}\medskip
\ifdefined\hypertarget\hypertarget{sourcepointer-1}{}\fi
\noindent\textit{Source reference.} Original lemma 00WI: sources/stacks-sites/pinned-original-source.tex, source lines 3541--3541. Navigation pointer invoked at proof line 3541; the counted span contains the pointer only and the cited result is not reproduced inside it, so the pointer is registered as a non-load-bearing navigation stub in the established form for Stacks material. No bibliography citation is involved.\par\medskip
\ifdefined\hypertarget\hypertarget{sourcepointer-2}{}\fi
\noindent\textit{Source reference.} Original lemma 00WT: sources/stacks-sites/pinned-original-source.tex, source lines 3556--3556. Navigation pointer invoked at proof line 3556; the counted span contains the pointer only and the cited result is not reproduced inside it, so the pointer is registered as a non-load-bearing navigation stub in the established form for Stacks material. No bibliography citation is involved.\par\medskip
\begin{lemma}
\label{lemma-quasi-compact}
Let $\mathcal{C}$ be a site. Let $U$ be an object of $\mathcal{C}$.
The following are equivalent
\begin{enumerate}
\item $U$ is quasi-compact, and
\item for every surjection of sheaves
$\coprod_{i \in I} \mathcal{F}_i \to h_U^\#$
there is a finite subset $J \subset I$ such that
$\coprod_{i \in J} \mathcal{F}_i \to h_U^\#$ is surjective.
\end{enumerate}
\end{lemma}

\begin{proof}
Assume (1) and let $\coprod_{i \in I} \mathcal{F}_i \to h_U^\#$
be a surjection. Then $\text{id}_U$ is a section of
$h_U^\#$ over $U$. Hence there exists a covering
$\{U_a \to U\}_{a \in A}$ and for each $a \in A$
a section $s_a$ of $\coprod_{i \in I} \mathcal{F}_i$
over $U_a$ mapping to $\text{id}_U$. By the construction of coproducts as
sheafification of coproducts of presheaves
(Lemma \ref{lemma-colimit-sheaves}), for each $a$
there exists a covering $\{U_{ab} \to U_a\}_{b \in B_a}$ and
for all $b \in B_a$ an $\iota(b) \in I$ and a section
$s_{b}$ of $\mathcal{F}_{\iota(b)}$ over $U_{ab}$
mapping to $\text{id}_U|_{U_{ab}}$. Thus after replacing
the covering $\{U_a \to U\}_{a \in A}$ by
$\{U_{ab} \to U\}_{a \in A, b \in B_a}$
we may assume we have a map $\iota : A \to I$
and for each $a \in A$ a section $s_a$ of $\mathcal{F}_{\iota(a)}$
over $U_a$ mapping to $\text{id}_U$.
Since $U$ is quasi-compact, there is a covering
$\{V_c \to U\}_{c \in C}$, a map $\alpha : C \to A$
with finite image, and $V_c \to U_{\alpha(c)}$ over $U$.
Then we see that $J = \Im(\iota \circ \alpha) \subset I$ works
because $\coprod_{c \in C} h_{V_c}^\# \to h_U^\#$ is surjective
(Lemma \ref{lemma-covering-surjective-after-sheafification})
and factors through $\coprod_{i \in J} \mathcal{F}_i \to h_U^\#$.
(Here we use that the composition
$h_{V_c}^\# \to h_{U_{\alpha(c)}}
\xrightarrow{s_{\alpha(c)}} \mathcal{F}_{\iota(\alpha(c))} \to h_U^\#$
is the map $h_{V_c}^\# \to h_U^\#$ coming from the morphism
$V_c \to U$ because $s_{\alpha(c)}$ maps to $\text{id}_U|_{U_{\alpha(c)}}$.)

\medskip\noindent
Assume (2). Let $\{U_i \to U\}_{i \in I}$ be a covering.
By Lemma \ref{lemma-covering-surjective-after-sheafification}
we see that $\coprod_{i \in I} h_{U_i}^\# \to h_U^\#$ is surjective.
Thus we find a finite subset $J \subset I$ such that
$\coprod_{j \in J} h_{U_j}^\# \to h_U^\#$ is surjective.
Then arguing as above we find a covering
$\{V_c \to U\}_{c \in C}$ of $U$ in $\mathcal{C}$
and a map $\iota : C \to J$
such that $\text{id}_U$ lifts to a section
of $s_c$ of $h_{U_{\iota(c)}}^\#$ over $V_c$.
Refining the covering even further we may assume
$s_c \in h_{U_{\iota(c)}}(V_c)$ mapping to $\text{id}_U$.
Then $s_c : V_c \to U_{\iota(c)}$ is a morphism over $U$
and we conclude.
\end{proof}

\end{document}
